\scanpage{12}
\deleted{
fixe, mettons $C$ \unsure{ici)} sont déterminés à ext. alg. près. Nous \unsure{allons} examiner la

\uline{Conjecture} \inserted{\uline{de transcendance}} Si $k$ est \unsure{pré-minimal|pré-maximal} pour $M$, alors
\[
\operatorname{deg\ tr}(k(g_0))=\dim G^{\unsure{u|n}}\qquad(=\dim G-1),
\]
$\unsure{\xi_0|g_0}$ \inserted{(\unsure{localité de $g_0$})} est un point \uline{maximal} de $G^{\unsure{u|n}}$ (du moins si $G$ et $G^{\unsure{u|n}}$ connexes).

\uline{Corollaire}. \inserted{Si $G$ connexe} $G^{\unsure{u|n}}$ est le $k$-groupe algébrique engendré par $g_0$. Donc $G$ est le $k$-groupe algébrique engendré par $g_0$ et $j(G_m)$ (les homothéties). Enfin $G$ est engendré sur $k$ par $g$ si $G^{\unsure{u|n}}$ est $\frac12$ simple et $\pi$ non algébrique sur $k$ $\deleted{\unsure{\longleftrightarrow}}$ ...

Contient le th. de Schneider, pour les courbes elliptiques sans mult. complexe. En fait, il signifie que le degré de transc. du système
\diagram{\begin{array}{cccc}
(\unsure{\zeta|\xi}^{-1}\omega_1,&\unsure{\zeta|\xi}^{-1}\omega_2,&\unsure{\zeta|\xi}^{-1}\eta_1,&\unsure{\zeta|\xi}^{-1}\eta_2)\\
\text{''}\ \omega'_1&\text{''}\ \omega'_2&\text{''}\ \eta'_1&\text{''}\ \eta'_2
\end{array}}
sur $Q$ est $3$ \deleted{\unsure{si}} (la seule relation étant
\[
\omega'_1\eta'_2-\omega'_2\eta'_1=1,
\]
i.e. l'équation de $G^{\unsure{u|n}}=\mathrm{Sl}(2)\subset G=\mathrm{Gl}(\unsure{2|7})$) pour une courbe elliptique sans mult. complexe définie sur un corps de nbs alg. (disons sur $\overline Q$).
Si l'invariant $j$ est transcendant, on trouve que $(j,\omega'_1,\omega'_2,\eta'_1,\eta'_2)$ ont le degré de \uline{trans. $4$} sur $Q$.
}
\margin{\unsure{connu!} pour courbes elliptiques à invariant transcendant et $\tau$ algébrique.....\unsure{?!}}
\margin{et cas \unsure{général.} \deleted{\unsure{$\beta s W$}} \deleted{$\pi$ non Alg. sur $k$?}}
\mspage{11}
